% ============================================
% ARTICLE VI: EVENTS, ACTIVITIES & PROJECTS
% ============================================
\article{Events, Activities \& Projects}

\gsection{Classification of Activities and Events}

For the purposes of this document, the following definitions apply:

\begin{enumerate}[label=\alph*), leftmargin=2em, itemsep=0.3em]
    \item \textbf{Activities:} Ongoing or periodic Club operations that occur throughout the academic year. Activities include membership functions, administrative processes, and continuous initiatives.
    \item \textbf{Events:} Scheduled, time-bound occurrences organized by the Club that involve planning, execution, and conclusion. Events may be internal (members-only) or external (open to the college community or public).
\end{enumerate}

\gsection{Annual Program of Activities and Events}

In alignment with its mission, the Club shall organize the following each academic year:
\\
\textbf{Mandatory Technical Events:}
\begin{enumerate}[label=\alph*), leftmargin=2em, itemsep=0.3em]
    \item \textbf{Bootcamps:} Minimum 2 per year --- Intensive training sessions on foundational technologies (e.g., Python, web development, mobile development).
    \item \textbf{Workshops:} Minimum 2 per year --- Hands-on learning sessions focused on specific technical skills or tools.
    \item \textbf{Competitions:} Minimum 2 per year --- Programming contests or technical challenges designed to foster problem-solving abilities.
    \item \textbf{Seminars/Tech Talks:} Minimum 1 per year --- Sessions with industry professionals, alumni, or academic experts.
\end{enumerate}
\\
\textbf{Club Activities:}
\begin{enumerate}[label=\alph*), leftmargin=2em, itemsep=0.3em]
    \item \textbf{Recruitment Drive:} Annual process for selecting new members, conducted in accordance with Article II.
    \item \textbf{Reformation:} Annual election and leadership transition process, conducted in accordance with Article VIII.
    \item \textbf{Orientation/Introduction:} Welcome session for new students at the beginning of each academic year.
    \item \textbf{Welcome and Farewell Program:} A formal program conducted after recruitment to welcome newly selected members and acknowledge departing members at the end of a tenure.
    \item \textbf{Club Project(s):} At least one practical project developed by Club members annually. The Club may undertake up to 2 projects simultaneously if resources permit.
    \item \textbf{Annual Report:} Year-end summary of Club activities prepared for the Department.
\end{enumerate}
\\
\textbf{Optional Activities:}
\begin{enumerate}[label=\alph*), leftmargin=2em, itemsep=0.3em]
    \item \textbf{Hackathons:} Problem-solving events where students build projects. These require significant planning and may be organized when the Club has sufficient capacity and member interest.
    \item \textbf{Webinars:} Online seminar sessions that may supplement or replace in-person seminars when guest speakers are unavailable for physical attendance.
    \item \textbf{Social Activities:} The Club may organize social activities such as hikes, picnics, or team-building events for members when resources permit. These activities are not mandatory but are encouraged to build camaraderie.
\end{enumerate}

\gsection{Event Management}

\textbf{Planning Phase:}
\begin{enumerate}[label=\alph*), leftmargin=2em, itemsep=0.3em]
    \item For each major event, the Event Manager shall develop a comprehensive plan including:
    \begin{itemize}[leftmargin=1.5em, itemsep=0.2em]
        \item Event objectives and expected outcomes
        \item Target audience and expected attendance
        \item Date, time, and venue selection with justification
        \item Detailed logistics requirements
        \item Promotional plan and timeline
        \item Preliminary budget requirements (prepared in coordination with the Treasurer)
    \end{itemize}
    \item The event plan requires approval from the Executive Board before proceeding.
    \item For events requiring guest speakers or external trainers, the Public Relations Officer shall initiate invitations at least 3 weeks in advance.
\end{enumerate}
\\
\textbf{Execution Phase:}
\begin{enumerate}[label=\alph*), leftmargin=2em, itemsep=0.3em]
    \item The Event Manager, with the support of all members, is responsible for the successful execution of the event.
    \item Members shall be assigned specific duties by the Event Manager based on the operational support roles outlined in Article II.
    \item The President or Secretary shall be present at all major events to provide Executive oversight.
    \item Any deviation from the approved plan costing more than 10\% of the budget requires immediate approval from the President.
\end{enumerate}
\\
\textbf{Post-Event Phase:}
\begin{enumerate}[label=\alph*), leftmargin=2em, itemsep=0.3em]
    \item After each event, the Event Manager shall submit a report to the President documenting:
    \begin{itemize}[leftmargin=1.5em, itemsep=0.2em]
        \item Event outcomes against objectives
        \item Member participation records
        \item Budget compliance summary
        \item Challenges encountered and lessons learned
        \item Recommendations for future similar events
    \end{itemize}
    \item This report, along with the Event Debrief Meeting, shall serve as a reference for future planning.
    \item The Event Manager shall ensure all borrowed equipment is returned and the venue is restored to its original condition.
\end{enumerate}
\\
\textbf{Multi-Day Event Procedures:}
\begin{enumerate}[label=\alph*), leftmargin=2em, itemsep=0.3em]
    \item For multi-day events, informal daily reviews may be conducted to address immediate operational needs.
    \item A formal debrief shall be held after the event's conclusion, as specified in Article IV.
    \item Responsibilities may be rotated across days, with coordination managed by the Event Manager.
    \item All members are expected to be present on the final day unless they have received prior approval for absence from the President.
\end{enumerate}

\gsection{Activity Periods}

\begin{enumerate}[label=\alph*), leftmargin=2em, itemsep=0.3em]
    \item The Club's main activities shall be scheduled during the regular academic year.
    \item All mandatory activities will be suspended during official college examination periods to allow members to focus on their academic responsibilities.
    \item Event planning and preparation may continue during examination periods but shall not require member attendance at meetings or events.
    \item The Executive Board shall publish an annual activity calendar at the beginning of each academic year.
\end{enumerate}

\gsection{Club Projects}

\textbf{Project Requirements:}
\begin{enumerate}[label=\alph*), leftmargin=2em, itemsep=0.3em]
    \item Each year, the Club shall undertake at least one project through the Project Development Unit (PDU) that:
    \begin{itemize}[leftmargin=1.5em, itemsep=0.2em]
        \item Addresses a real need for the campus, Club, or community
        \item Involves Club members in actual development work
        \item Demonstrates the Club's technical capabilities
        \item Provides practical experience to participating members
    \end{itemize}
    \item Projects may include applications for campus use, tools for the Club itself, or solutions for local community needs.
    \item The PDU may also prepare club-led submissions, prototypes, and demonstrations for hackathons, ideathons, competitions, and external project showcases when suitable opportunities arise.
\end{enumerate}
\\
\textbf{Team Composition:}
\begin{enumerate}[label=\alph*), leftmargin=2em, itemsep=0.3em]
    \item Internal project teams shall consist of 4 members maximum, including the Project Lead.
    \item This limit ensures active contribution from all team members and prevents free-riding.
    \item The Project Lead, acting through the PDU, has the authority to select team members based on skills, interest, and availability.
    \item A minimum of 3 members (including the Project Lead) is required for a project team to be formally recognized.
\end{enumerate}
\\
\textbf{Project Management:}
\begin{enumerate}[label=\alph*), leftmargin=2em, itemsep=0.3em]
    \item The Project Lead shall establish milestones, assign tasks, and monitor progress through the PDU.
    \item Regular progress updates shall be provided to the Executive Board at general meetings.
    \item The Project Lead has \textbf{Bench Authority} as specified in Article III to temporarily sideline non-performing team members.
    \item Upon project completion, the Project Lead shall submit a final report including technical documentation, member contributions, and outcomes to the President.
\end{enumerate}

\gsection{Project Development Unit Readiness and External Events}

\begin{enumerate}[label=\alph*), leftmargin=2em, itemsep=0.3em]
    \item The PDU shall remain ready to support external technical opportunities throughout the academic year, including hackathons, ideathons, competitions, project demos, exhibitions, and similar events.
    \item The Project Lead shall maintain a prepared list of members who can be mobilized quickly for external event participation based on skill, availability, and event requirements.
    \item Participation in external events shall be coordinated with the Executive Board, and the President shall have final approval over official Club representation.
    \item During official college examination periods, the PDU's external event obligations shall be suspended unless the Executive Board determines that participation is urgent and approves it in advance.
    \item The PDU may organize internal preparation sessions, prototype reviews, and demo rehearsals to ensure the Club is ready for external opportunities when they arise.
\end{enumerate}

\gsection{Participation in YATRA}

The Club plays an important role in YATRA, the college's national-level event, working alongside other clubs and the Free Student Council (FSC). The Executive Board shall coordinate with the YATRA organizing committee to determine the Club's specific contributions each year.
